\documentclass[Markos_Delaportas_1316.tex]{subfiles}

\begin{document}

\section*{Introduction}
\addcontentsline{toc}{section}{\protect\numberline{}Introduction}
Over the Air Computing offers a paradigm shift in how distributed 
computations are performed in wireless networks. By exploiting the 
inherent properties of the wireless medium, \gls{ota} enables simultaneous 
transmissions from multiple devices to be combined in a way that directly 
computes a desired function of the transmitted data. This technique has 
garnered significant attention due to its potential to enhance efficiency 
and reduce latency in applications such as federated learning, distributed 
sensing, and real-time analytics. The motivation behind this thesis is to 
explore the practical implementation of \gls{ota} using Software Defined 
Radios \gls{sdr} and MATLAB, addressing the challenges associated with 
synchronization, channel impairments, and waveform design. The research 
objectives include designing a prototype \gls{ota} network, validating its 
performance, and analyzing the impact of various impairments on computation 
accuracy.
\subsection{Background and Motivation}
The increasing need for efficiency in computation-oriented applications
where the interest lies in a function of local information, rather than
the information itself \cite{sahin_survey_2023}. The conventional
approach of separating communication and computation often proves
suboptimal \cite{sahin_survey_2023,goldenbaum_analyzing_2011,
goldenbaum_robust_2013}. The potential of \gls{ota} Computation
is to achieve significantly higher computation rates by harnessing
interference through signal superposition in a wireless Multiple Access
Channel (MAC) \cite{sahin_survey_2023}.
\subsection{Problem Statement and Research Objectives}
Focusing on designing and validating a practical \gls{ota} network layout
suitable for implementation in a controllable
environment\cite{sahin_survey_2023}. The specific role of Software
Defined Radios (SDRs) and MATLAB/simulation tools for prototype
validation and addressing practical limitations such as synchronization
and channel impairments \cite{sahin_survey_2023,xiao_heterogeneous_2025,
evgenidis_waveform_2025}.
\subsection{Thesis Contributions and Organization}

    \ifSubfilesClassLoaded{
        \clearpage
        \printbibliography
    }{}

\end{document}
